\scanpage{69}
exacte
\[
0\mathbin{-}\delta'\mathbin{-}\delta\mathbin{-}\delta''\mathbin{-}0
\]
et la suite exacte associée
\[
0\mathbin{-}{}_n\delta'\mathbin{-}{}_n\delta\mathbin{-}{}_n\delta''\mathbin{-}\delta'_n\mathbin{-}\delta_n\mathbin{-}\delta''_n
\]
Reste : prouver : \textcircled{\unsure{a|0}}) le fait que $\Gamma$ est sans torsion. Mais $\Gamma=\mathbf Z/d\mathbf Z$, \unsure{où|i.e.|:=} $d$ est le pgcd des $d_i$, $X_0=\sum d_iD_i$, \& $X_{0\mathrm{red}}=\sum\unsure{D_i}$. \deleted{(On \unsure{voit|sait} que si $\exists$ sur $X$ un $0$-cycle \unsure{un} $0$-cycle de degré $1$} si on a un pt fermé $x$ de $X_\eta$, son adhérence \inserted{schématique dans $X$} est un sous-schéma \inserted{fermé $Z$} \deleted{plat} de $X$, plat et fini sur $S$, \deleted{qui est} dont le degré sur $S$ est aussi égal à la multiplicité d'intersection \inserted{$Z.X_0$} sur $X$, qui est $\sum d_iZ.D_i$, donc multiple de $d$. Donc si $d'\deleted{\unsure{d}}$ est le pgcd des degrés des 0-cycles sur $X_\eta$, on a $d\mid d'$. Comme $d'=1$, on a $d=1$.
\margin{donner \uline{lemme} (ici \unsure{sans|sur} $X$ régulier, semble-t-il)}

\textcircled{b} $\Gamma\longrightarrow\operatorname{NS}(X_0)$ est \uline{injectif}.
\margin{\uline{Donner cor. Néron} \uline{ou th.}}
[Soit $R$ son noyau, \inserted{qui est $\mathbf Z$-libre de t.f. \unsure{p. $0$|$r=0$|$\ne0$})}, \deleted{\unsure{On}} a des suites exactes
\[
0\mathbin{-}\mathrm H^1(X,\mu_n)\mathbin{-}\mathrm H^1(X_{\bar\eta},\mu_n)^I\unsure{\mathbin{-}|\longrightarrow}{}_{\unsure{n|a}}\Gamma_n\mathbin{-}\operatorname{NS}(X_0)_n
\]
d'où, par passage à la limite sur les puissances d'un $\ell$
\[
0\mathbin{-}\mathrm H^1(X,\mathbf Z_\ell(1))\mathbin{-}\mathrm H^1(X_{\bar\eta},\mathbf Z_\ell(1))^I\longrightarrow\deleted{T_\ell}(\Gamma\deleted{\unsure{)}}\otimes\mathbf Z_\ell\mathbin{-}\deleted{T_\ell}(\operatorname{NS}(X_0)\deleted{\unsure{)}}\otimes\mathbf Z_\ell
\]
i.e. on trouve, comme $\deleted{T_\ell(R)}\unsure{\deleted{\simeq}|\simeq}R\otimes\mathbf Z_\ell\simeq\operatorname{Ker}(\Gamma\otimes\mathbf Z_\ell\mathbin{-}\operatorname{NS}(X_0)\otimes\mathbf Z_\ell)$
\deleted{$T_\ell(R$}\quad$\unsure{\beta|\rho}$
\[
0\mathbin{-}\mathrm H^1(X,\mathbf Z_\ell(1))\mathbin{-}\mathrm H^1(X_{\bar\eta},\mathbf Z_\ell(1))^I\mathbin{-}R\otimes\mathbf Z_\ell\mathbin{-}0
\]
Donc notre assertion équivaut à dire que l'on a la formule des cycles invariants ; \deleted{\unsure{cette|c'est}} il suffit \unsure{$\subset$|c.} de la prouver à torsion près.]
\margin{[\quad] en remarque après le cor.}

\deleted{L'idée est de se ramener au cas où $X$ est de dim relative un. Dans ce cas, on \uline{sait} que $\Gamma$ est nul (c'est le fait que une certaine matrice \inserted{d'intersection $(D_i.D_j)$} est de rang $\deleted{\unsure{\ell|r|(}}r-1)$, où $r$ est le nb de composantes $D_i$)}
Or c'est fait par Raynaud dans sa note : \unsure{s|S}pécialisation des foncteurs de Picard, \uline{th. 3} --- dans version sans \deleted{\unsure{dans|l'hyp.}} hyp. projective.
\margin{donner dém. directe \unsure{au|en} cas de dim. rel. 1}
\deleted{[Il \unsure{faut} \unsure{que} \unsure{l'on} \unsure{suppose} que $R^{\unsure{1|I}}_{\unsure{X,\Lambda}}$ \unsure{est fini} \unsure{si} $\unsure{d}=\unsure{f|p}$,
\unsure{des} \illegible{1} \unsure{comme} \unsure{cas} \unsure{du|de} \illegible{1} $\operatorname{NS}(X_0)$ \unsure{est nul}, \illegible{1} $\unsure{=}$ \illegible{2}\unsure{)}]}
\margin{68}
